\documentclass{article}
\usepackage[UTF8]{ctex}
\usepackage{amsmath,amssymb,amsfonts}
\usepackage{enumitem}
\usepackage{geometry}
\geometry{a4paper, margin=1in}
% === 错题本排版优化（不参与题目标识） ===
\setlength{\parindent}{0pt}
\setlength{\parskip}{0.35em plus 0.10em minus 0.08em}
\setlength{\emergencystretch}{3em}
\setlist{itemsep=0.25em, topsep=0.40em, parsep=0pt}
\allowdisplaybreaks[2]
% === 错题本排版优化结束 ===
\begin{document}
\section*{重庆统考2026年高等代数真题}
\section*{54．重庆统考2026年高等代数试卷}
一．计算题．\\

\subsection{题目1}
1．求多项式 \(f(x)=4 x^{4}-2 x^{3}-16 x^{2}+5 x+9\) 与 \(g(x)=2 x^{3}-x^{2}-5 x+4\) 的公共根．

\subsection{题目2}
2．计算行列式

\[
D_{n}=\left|\begin{array}{cccc}
1+a b & 1+a b^{2} & \cdots & 1+a b^{n} \\
1+a^{2} b & 1+a^{2} b^{2} & \cdots & 1+a^{2} b^{n} \\
\vdots & \vdots & & \vdots \\
1+a^{n} b & 1+a^{n} b^{2} & \cdots & 1+a^{n} b^{n}
\end{array}\right| .
\]

\subsection{题目3}
3．已知实二次型

\[
f\left(x_{1}, x_{2}, x_{3}\right)=2 x_{1}^{2}+5 x_{2}^{2}+5 x_{3}^{2}+2 a x_{1} x_{2}+2 b x_{1} x_{3}-8 x_{2} x_{3}(a>0)
\]

在正交变换下可化为 \(y_{1}^{2}+y_{2}^{2}+c y_{3}^{2}\) 。

\begin{itemize}
  \item[（1）] 写出二次型对应的实对称矩阵 \(A\) ，给出 \(c\) 的值．
  \item[（2）] 求矩阵 \(A-E\) 的秩，并求 \(a, b\) 。
\end{itemize}

\subsection{题目4}
4．已知 \(A\) 为3阶实对称矩阵，其特征值为 \(\lambda_{1}=1, \lambda_{2}=-1, \lambda_{3}=0, \alpha_{1}=(1,2,2)^{\mathrm{T}}, \alpha_{2}=(2,1,-2)^{\mathrm{T}}\)分别是特征值 \(\lambda_{1}, \lambda_{2}\) 对应的特征向量，求矩阵 \(A\) 。

\subsection{题目5}
5．设 \(P^{2 \times 2}\) 中的矩阵

\[
G_{1}=\left(\begin{array}{ll}
a & 1 \\
1 & 1
\end{array}\right), G_{2}=\left(\begin{array}{cc}
1 & a \\
1 & 1
\end{array}\right), G_{3}=\left(\begin{array}{cc}
1 & 1 \\
a & 1
\end{array}\right), G_{4}=\left(\begin{array}{cc}
1 & 1 \\
1 & a
\end{array}\right) .
\]

讨论由 \(G_{1}, G_{2}, G_{3}, G_{4}\) 生成的子空间 \(L\left(G_{1}, G_{2}, G_{3}, G_{4}\right)\) 的维数并给出相应的一组基。

\subsection{题目6}
6．线性空间 \(P^{3}\) 上的线性变换 \(\sigma\) 在基 \(\alpha_{1}, \alpha_{2}, \alpha_{3}\) 下的矩阵为

\[
A=\left(\begin{array}{ccc}
1 & -1 & 1 \\
1 & 3 & -1 \\
1 & 1 & 1
\end{array}\right)
\]

向量 \(\xi\) 在 \(\alpha_{1}, \alpha_{2}, \alpha_{3}\) 下的坐标为 \((1,1,1)^{\mathrm{T}}\) ．

\begin{itemize}
  \item[（1）] 求 \(\sigma(\xi)\) 。
  \item[（2）] 求 \(\sigma\) 的特征值和特征向量。
  \item[（3）] 判断是否存在 \(P^{3}\) 的一组基，使得 \(\sigma\) 在这组基下的矩阵为对角矩阵。
\end{itemize}

\subsection{题目7}
7．给定线性方程组

\[
\left\{\begin{array}{l}
x_{1}+x_{2}+x_{3}+x_{4}+x_{5}=1 \\
-x_{1}+x_{3}+3 x_{4}+x_{5}=a \\
x_{1}+2 x_{2}+3 x_{3}+5 x_{4}+3 x_{5}=3 \\
3 x_{1}+2 x_{2}+x_{3}-x_{4}+x_{5}=b
\end{array}\right.
\]

\begin{itemize}
  \item[（1）] 当 \(a, b\) 取何值时，方程组有解？
  \item[（2）] 当方程组有解时，求方程组的所有解。
  \item[（3）] 当方程组有解时，求所有解的极大线性无关组。
\end{itemize}

\subsection{题目8}
8．已知 \(A\) 为数域 \(P\) 上的 6 阶矩阵，\(f(\lambda)=(\lambda-2)^{3}(\lambda+3)^{2}(\lambda-4)\) 为 \(A\) 的特征多项式，\(m(\lambda)=\) \((\lambda-2)^{2}(\lambda+3)(\lambda-4)\) 为 \(A\) 的最小多项式．

\begin{itemize}
  \item[（1）] 求 \(A\) 的所有不变因子．
  \item[（2）] 写出 \(A\) 的 Jordan 标准形．
  \item[（3）] 写出 \(A\) 的有理标准形．
\end{itemize}
\end{document}
